\section{Full relative asymptotics for coprime circles}
\label{sec:relativecoprime}
The growing-entry comparison closes the arithmetic count at
relative precision. This is stronger than an exponential-base
estimate and treats every modulus, with the parity effect explicit.

\begin{lemma}[An elementary sufficiently uniform totient bound]
\label{lem:totientdiffusion}
For all \(n\geq3\),
\[
 n/\varphi(n)\ll\log\log(3n).
\]
\end{lemma}
\begin{proof}
We recall an elementary argument, not claiming this classical
estimate as new.
Chebyshev's upper bound \(\vartheta(x)\ll x\) and partial
summation give
\(\sum_{p\leq x}(\log p)/p\ll1+\log x\).
Take \(s=1+1/\log x\).
Expanding the two logarithmic Euler products, and using
\(1-e^{-u}\leq u\), gives
\[
 \log\prod_{p\leq x}
       \frac{1-p^{-s}}{1-p^{-1}}
 \leq(s-1)\sum_{p\leq x}\frac{\log p}{p-1}\ll1 .
\]
Thus
\[
 \prod_{p\leq x}(1-1/p)^{-1}
                  \ll\zeta(s)\ll\log x,
\]
using \(\zeta(s)\leq1+1/(s-1)\).
Set \(y=\log(3n)\).
There are at most \((\log n)/\log y\) prime divisors of
\(n\) exceeding \(y\), whose contribution to the logarithm
of the product is at most
\(2\log n/(y\log y)=O(1)\).
The small-prime contribution is \(O(\log y)\).
This proves the bound for large \(n\); a fixed constant
covers the remaining finite range.
No prime number theorem or RH hypothesis is used.
\end{proof}

\paragraph{Maintained full-domain density input.}
Lemma~\ref{lem:totientdiffusion} is now formalized for every
integer \(n\geq3\), with one positive absolute constant.
\texttt{HamiltonianMertensProduct} consumes the existing
cached classical Mertens product asymptotic and derives a
positive reciprocal bound.
\texttt{HamiltonianTotientEuler} proves the actual totient
Euler product and applies the previously proved large-prime
product estimate to the same prime divisors.
\texttt{HamiltonianTotientCutoff} retains both supported
products and proves, for every real \(y\geq4\),
\[
 \frac n{\varphi(n)}
 \leq
 \frac{\exp(2\log n/(y\log y))}
      {\prod_{p\leq y}(1-1/p)}.
\]
\texttt{HamiltonianTotientDiffusion} chooses
\(y=\log(3n)\). Above the fixed Mertens and logarithmic
cutoffs its tail factor is at most \(e^2\).
One constant covers all smaller \(n\geq3\) using
\(\log\log(3n)\geq\log\log3>0\) and \(n/\varphi(n)\leq n\).
The maintained theorem has the exact quantifiers
\(\exists C>0\ \forall n\geq3\).
The Mertens input is unconditional and its new consumer has
passed a foundational axiom audit; no dependency rebuild or
numerical range test is used.
This closes the density lemma.

\paragraph{Maintained quantitative odd-modulus count.}
\texttt{HamiltonianTotientSublog} proves, for every \(b>0\),
one threshold above which
\((n/\varphi(n))^2\leq b\log n\), uniformly over every modulus.
For every fixed \(K>0\), \texttt{HamiltonianCoprimeCountBudget}
supplies the simultaneous entry cap \(n/\varphi(n)\leq\log n\)
and the full loss budget
\(K(n/\varphi(n))^2\leq(\log n)/90\).
Thus the complete error
\(17n^{-1/9}\exp(K(n/\varphi(n))^2)\)
is at most \(17n^{-1/10}\).

\texttt{HamiltonianCoprimeAperiodicCount} applies these bounds
and the proved physical half-gap to the actual native
coprime kernel \(P_n=A_{X_n}/\varphi(n)\).
Writing \(B_n=P_n-J_n/n\) and
\[
 R_n=\frac{(n-1)!}{n^n}\det(I-B_n)
                      \det(I-B_nB_n^*)^{-1/2},
\]
the compiled theorem states that one \(N_0\geq256\) satisfies,
for every odd \(n\geq N_0\),
\[
 R_n>0,\qquad
 \left|\frac{Z_{\rm dir}(P_n)}{R_n}-1\right|
                      \leq17n^{-1/10}.
\]
Here \(Z_{\rm dir}\) is the actual single-cycle permutation sum.
The norm transfer between the native residue vertices and
the finite counting coordinates uses the proved norm-preserving
reindexing theorem explicitly.
Every new declaration passes the focused foundational audit.
This closes the actual odd stochastic count, including its
rate and positive reference.

\paragraph{Maintained full odd arithmetic normalization.}
\texttt{HamiltonianCoprimeConstant} proves positivity on
every supported divisor and identifies the paper's exact
real-power product \(\mathfrak C(n)\) with the positive
square root of the determinant quotient. It is therefore
the actual determinant--Gram factor in \(R_n\).
\texttt{HamiltonianCoprimeUnorientedCount} restores the
edge weight and the reversal factor exactly.
Writing
\[
 H_n=\frac{(n-1)!}{2n^n}\varphi(n)^n\mathfrak C(n),
\]
its compiled theorem gives one \(N_0\geq256\) such that,
for every odd \(n\geq N_0\),
\[
 H_n>0,\qquad
 \left|\frac{h(X_n)}{H_n}-1\right|\leq17n^{-1/10}.
\]
This is the displayed factorial normalization of the actual
unoriented arithmetic graph, with no restriction on its
distinct prime factors or prime-power exponents.
All fourteen new declarations pass the focused foundational
audit. The constant's infinite-series analysis and uniform
bounds, the final Stirling normalization and the even-modulus
count remain formalization obligations.

\paragraph{Maintained exact arithmetic determinant.}
\texttt{HamiltonianFourierDeterminant} proves that the determinant
of the complete cyclic difference operator is the product of its
full Fourier symbol. Both frame determinants cancel the exact
grid normalization.
\texttt{HamiltonianCoprimeDeterminant} applies this to the
actual compensated coprime kernel, for every positive modulus
and every real scalar \(z\):
\[
 \det(I-zB_n)
   =\prod_{t\mid n}F_z(t)^{\varphi(t)},\qquad
 F_z(1)=1,\qquad
 F_z(t)=1-z\frac{\mu(t)}{\varphi(t)}\quad(t>1).
\]
The constant mode is compensated before the determinant is taken.
The same formula retains the parity mode at even modulus and all
nonsquarefree divisors. The frequency product is regrouped using
the proved exact conductor multiplicities on the actual divisor
support. All ten new declarations pass the focused foundational
audit. Thus both determinants in the symmetric quotient are
explicit finite arithmetic products. The positive square-root
identification with the displayed \(\mathfrak C(n)\) and the
actual odd unoriented factorial count are now compiled as
recorded above. Its infinite series and bounds, final Stirling
normalization and the even count remain separate obligations.


For odd \(n\geq3\), define the positive arithmetic factor
\[
 {\mathfrak C}(n)=
 \prod_{\substack{t\mid n\\t>1}}
 \left(\frac{\varphi(t)-\mu(t)}
            {\varphi(t)+\mu(t)}\right)^{\varphi(t)/2}.
\]
Every denominator is positive, since an odd \(t>1\) has
\(\varphi(t)\geq2\).
Factors with \(\mu(t)=0\) equal one.

\begin{theorem}[Relative enumeration through every modulus]
\label{thm:relativecoprime}
As \(n\to\infty\), the number of unoriented coprime circles is
\[
 h(X_n)=
 \begin{cases}
 (1+O(n^{-1/10}))\,\pi\,(\varphi(n)/e)^n,&n\ {\rm even},\\[2mm]
 (1+O(n^{-1/10}))\,{\mathfrak C}(n)
       \sqrt{\dfrac{\pi}{2n}}\,(\varphi(n)/e)^n,
                                      &n\ {\rm odd}.
 \end{cases}
\]
The implied constant is absolute for sufficiently large \(n\).
The estimate holds along all moduli of the indicated parity,
without restrictions on their prime factors.
For odd moduli the equivalent factorial form is
\[
 h(X_n)=(1+O(n^{-1/10}))\,
       \frac{(n-1)!}{2n^n}\,\varphi(n)^n\,{\mathfrak C}(n).
\]
\end{theorem}
\begin{proof}
Use the exact spectrum proved in Theorem~\ref{thm:coprimecircles}.
The normalized adjacency \(P=A_{X_n}/\varphi(n)\) has a
fixed gap after removing the constant vector and, only in
the even case, its exact parity vector.
The largest entry is \(1/\varphi(n)\).
By Lemma~\ref{lem:totientdiffusion}, its dimension times this
entry is \(O(\log\log n)\).

If \(n=2m\), its cross block \(T\) is doubly stochastic,
has nontrivial singular norm at most \(1/2\), and satisfies
\[
 m\max T_{ij}\leq n/\varphi(n)=O(\log\log n)=o(\log m).
\]
Theorem~\ref{thm:relativebipartite} therefore gives
\(Z_{\rm H}(P)=(1+o(1))\pi e^{-n}\).
Every Hamilton cycle has weight \(\varphi(n)^{-n}\),
which proves the even formula.

If \(n\) is odd, \(B=P-J_n/n\) has norm at most \(1/2\), and
\(A=O(\log\log n)\) satisfies \(A^2=o(\log n)\).
Apply Theorem~\ref{thm:relativeaperiodic}.
For each divisor \(t>1\) of \(n\), the eigenvalue
\(\mu(t)/\varphi(t)\) has multiplicity \(\varphi(t)\).
The determinant ratio for symmetric \(B\) is exactly
\[
 \frac{\det(I-B)}{\sqrt{\det(I-B^2)}}
   =\prod_{\substack{t\mid n\\t>1}}
       \left(\frac{1-\mu(t)/\varphi(t)}
                  {1+\mu(t)/\varphi(t)}\right)^{\varphi(t)/2}
   ={\mathfrak C}(n).
\]
Divide the directed count by two and restore the common
edge-weight factor \(\varphi(n)^n\).
This proves the factorial formula, and Stirling proves the
displayed odd-order normalization.
Choose, for example, the common gap parameter \(\delta=1/3\).
The relative comparison errors are at most
\(O(n^{-1/9}\exp(C(\log\log(3n))^2))\), uniformly over
all moduli. Since \((\log\log(3n))^2=o(\log n)\), the
exponential factor is at most \(n^{1/90}\) eventually.
The permanent error \(O(n^{-1/8})\) and Stirling error
\(O(n^{-1})\) are smaller. Thus the asserted common
\(O(n^{-1/10})\) error follows, with an absolute constant;
no numerical threshold is asserted.
There is no unproved prime-sampling estimate in this consumer.
\end{proof}

\begin{proposition}[The prime-factor dependence of the odd constant]
\label{prop:coprimeconstant}
The factor \({\mathfrak C}(n)\) depends only on the distinct
prime divisors of odd \(n\), and
\[
 \log{\mathfrak C}(n)
 =1+\sum_{j\geq1}\frac{
       1-\prod_{p\mid n}(1-(p-1)^{-2j})}{2j+1}.
\]
The series converges absolutely and uniformly over all odd moduli.
In particular
\[
 e\leq{\mathfrak C}(n)\leq e^{11/9}.
\]
It increases when new distinct prime factors are added.
Along odd primorials it converges to the finite constant
obtained by taking the product over all odd primes in the
series above.
\end{proposition}
\begin{proof}
Only squarefree divisors contribute.
Expand \(\tfrac12\log((1-x)/(1+x))
=-\sum_{j\geq0}x^{2j+1}/(2j+1)\), valid for
\(|x|\leq1/2\).
The \(j=0\) sum is
\(-\sum_{t\mid n,t>1}\mu(t)=1\).
For \(j\geq1\), the multiplicative divisor sum is
\[
 \sum_{t\mid n}\frac{\mu(t)}{\varphi(t)^{2j}}
       =\prod_{p\mid n}(1-(p-1)^{-2j}),
\]
giving the displayed series.

Its numerator is nonnegative and at most
\[
 \sum_{\substack{p\ {\rm odd}\\p\ {\rm prime}}}(p-1)^{-2j}
   \leq\sum_{\ell\geq1}(2\ell)^{-2j}
   =4^{-j}\zeta(2j)\leq2\cdot4^{-j}.
\]
Here the elementary integral bound gives \(\zeta(2j)\leq2\).
This majorant proves uniform absolute convergence.
Moreover
\[
 0\leq\log{\mathfrak C}(n)-1
   \leq\frac23\sum_{j\geq1}4^{-j}=2/9.
\]
Adding a prime decreases each finite product inside the
series, and therefore increases the constant.
The same majorant justifies the primorial limit and the
absolutely convergent infinite-prime expression.
\end{proof}

\begin{remark}[Exact calibration families and what has been settled]
For odd prime powers \(n=p^a\),
\[
 {\mathfrak C}(n)
       =\left(\frac p{p-2}\right)^{(p-1)/2};
\]
in particular \({\mathfrak C}(3^a)=3\).
For prime \(n\), the formula agrees asymptotically with
the exact complete-graph count \((n-1)!/2\).
For powers of two, it agrees with the exact complete
bipartite count \((n/2)!(n/2-1)!/2\).
These are symbolic consistency checks on infinite families,
not finite experimental certificates.

The result determines the full relative asymptotic, including
its parity-dependent polynomial factor and arithmetic constant.
It does not merely count the already known constant-step
Hamilton cycles or reprove their existence.
The infinite-prime interaction appears in an explicit,
uniformly convergent expression for the odd constant.
Priority remains subject to primary-literature comparison.
\end{remark}
